\scanpage{270}
\mspage{82}
% archive sheet number: 270
$\{K,i,j,\bar f,\alpha\}$ est (sauf erreur) défini à iso. can. près par la condition de représenter le foncteur $\mathcal M\longmapsto\{F,G,u\}^{\unsure{\mathrm o|\circ}}(\mathcal M^{\mathrm o})$\ldots

En particulier, le système défini par $G\in\mathrm{Ob}\,\mathbb D_{\mathcal M}(J)$

\[
(f^*(\deleted{\unsure{F}}G),\deleted{\unsure{f|F}}G,\mathrm{id}:f^*(G)\longrightarrow f^*(G))
\]

est défini sauf erreur par $\bar f^*(G)$ sur $K(f)$. On se demande de trouver une interprétation similaire \deleted{\unsure{pour}}

\[
(F,f_*(F),f^*f_*F\longrightarrow F)
\]

($F\in\mathrm{Ob}\,\mathbb D(I)$). Le candidat naturel serait ici $\Phi=i_*(F)$. Admettons ce point (i.e. \inserted{surtout} la vérification du fait que $j^*(i_*(F))\simeq f_*(F)$).

Tout ceci montre que si $\mathbb D$ est à nouveau un dérivateur quelconque, on peut, par les formules précédentes, trouver un objet canonique \inserted{$\Phi$} de $\mathbb D(K)$, en termes soit d'un \inserted{$\mathbb D$-coeff} $F$ sur $I$ (en prenant $\Phi=i_{\unsure{1|*}}(F)$), soit d'un $\mathbb D$-coeff: $G$ sur $J$ (en prenant $\Phi=\bar f^*(\unsure{G})$.

D'autre part, on peut espérer définir canoniquement

\[
\begin{aligned}
\mathbb D(K)\deleted{\unsure{(f)}}&\longrightarrow\mathbb D(I\times\mathbb{\Delta}_1)\\
\mathbb D(K)&\longrightarrow\mathbb D(J\times\mathbb{\Delta}_1)
\end{aligned}
\]

dont les composés avec les foncteurs ``squelette''

\[
\begin{aligned}
\mathbb D(I\times\mathbb{\Delta}_1)&\longrightarrow\underline{\operatorname{Hom}}(\mathbb{\Delta}_1^{\mathrm o},\mathbb D(I))\\
\mathbb D(J\times\mathbb{\Delta}_1)&\unsure{\text{---}|\longrightarrow}\underline{\operatorname{Hom}}(\mathbb{\Delta}_1^{\mathrm o},\mathbb D(J))
\end{aligned}
\]
